% ============================================================
%  §5  Structural Interpretation
% ============================================================

\section{Structural Interpretation}
\label{sec:structural}

\subsection{Solution Classes of the Unified Equation}

The unified VED$\times$IFGT equation admits several classes of
solutions that correspond to distinct cognitive phenomena.
We identify the principal classes here; detailed analysis is
deferred to future work.

\begin{center}
\renewcommand{\arraystretch}{1.5}
\begin{tabular}{p{0.29\textwidth}p{0.31\textwidth}p{0.30\textwidth}}
\toprule
\textbf{Solution class} & \textbf{Field structure} & \textbf{Cognitive interpretation} \\
\midrule
Smooth gradient region   & $\nabla\Phi$ slowly varying     & Background cognitive disposition \\
Shallow metastable well  & $\nabla\Phi\approx 0$, small depth & Tentative or unstable concept \\
Deep stable well         & $\nabla\Phi\approx 0$, large depth & Entrenched concept or belief \\
Limit-cycle trajectory   & $\JI$ circulates without sedimentation & Recurring thought pattern \\
Barrier region           & $|\nabla\Phi|$ large             & Cognitive block or resistance \\
Ridge structure          & Local maximum of $\Phi$           & Conceptual boundary or category edge \\
\bottomrule
\end{tabular}
\end{center}

\subsection{Depth, Width, and Robustness}

Three geometric parameters characterize a conceptual attractor:

\textbf{Depth} ($-\min\Phi$ within the well) measures resistance to
deformation by new input. Deep concepts are stable under perturbation;
shallow concepts are easily displaced.

\textbf{Width} (spatial extent of the sub-threshold region) measures
contextual range. Broad concepts apply across many domains; narrow
concepts are context-specific.

\textbf{Robustness} (curvature of $\Phi$ at the minimum) measures
precision. High curvature produces sharp, well-defined concepts;
low curvature produces fuzzy or polysemous ones.

These parameters are not fixed intrinsic properties; they evolve
continuously with $\Phi$ under the dynamics of
Eqs.~\eqref{eq:I_evolution}--\eqref{eq:sigma_barrier}.

\subsection{Conceptual Boundaries}

Ridge structures in $\Phi$ --- regions of locally high potential
separating neighboring wells --- correspond to conceptual boundaries:
the distinctions between concepts.
These boundaries are not sharp lines but extended regions of high
traversal cost.

A key consequence is that conceptual boundaries shift under repeated
linguistic or experiential perturbation.
Tokens that consistently couple to the gradient near a ridge will
either reinforce it (deepening the distinction) or erode it (merging
the adjacent concepts).
This provides a mechanistic account of semantic drift and conceptual
change.

\subsection{The Vocabulary as a Distribution}

The full set of stabilized tokens $\{\ell_i\}$ constitutes a vocabulary.
In this framework, vocabulary is not a list but a distribution over
stability measures:
\begin{equation}
  \mathcal{V} = \bigl\{\ell_i \;\big|\; S(\ell_i) > \epsilon \bigr\},
  \label{eq:vocabulary}
\end{equation}
where $S(\ell_i)$ is a stability functional derived from reuse frequency,
persistence time, and basin depth, and $\epsilon$ is a system-dependent
threshold.

Vocabulary is a field property, continuously reshaped by use.
Token frequencies reflect basin accessibility, not semantic importance.
Boundaries between tokens are fuzzy and context-sensitive.
